% ============================================================
%  Retrieval Results  % ============================================================
%  Retrieval Results  % ============================================================
%  Retrieval Results  % ============================================================
%  Retrieval Results  \input{5_retrieval_results} in main paper
% ============================================================

\definecolor{sectiongray}{gray}{0.95}

\section{Results}
\label{sec:results}

We address four questions:
\textbf{RQ1} : does query-conditioned source selection improve retrieval over a single-index baseline?
\textbf{RQ2} : are query routing (QRe) and uniform multi-source allocation (UMS) complementary remedies for size-imbalanced retrieval?
\textbf{RQ3} : does query-conditioned chunking optimisation improve retrieval over a global baseline?
\textbf{RQ4} : does query-conditioned \emph{metadata indexing} (filters and rerankers on \textit{Les Débats}) improve retrieval over a global baseline?

\subsection{Retrieval Strategies}
\label{sec:retrieval-results}

We evaluate the four retrieval strategies of
Section~\ref{sec:retrieval-strategies} on the multi-hop subset of
HistoriQA-ThirdRepublic~\cite{historiqa-pellet}, broken down by question type:
\textit{cross-newspaper} (314 q., both gold documents from newspapers) and
\textit{newspaper $\rightarrow$ Débats} (571 q., one newspaper plus one
parliamentary document). We include comparison to BM25~\cite{bm25} for sparse retrieval and 2 graph-based external baseline: HippoRAGv2~\cite{hipporagv2} and LinearRAG~\cite{linearrag_2025}.


\begin{table*}[ht]
\centering
\small
\setlength{\tabcolsep}{4pt}
\begin{tabular*}{\textwidth}{@{\extracolsep{\fill}} l ccc ccc ccc @{}}
\toprule
 & \multicolumn{3}{c}{\textbf{Overall}} & \multicolumn{3}{c}{\textbf{Cross-newspaper} (314 q.)} & \multicolumn{3}{c}{\textbf{Newspaper-Débats} (571 q.)} \\
\cmidrule(lr){2-4} \cmidrule(lr){5-7} \cmidrule(lr){8-10}
\textbf{Configuration} & R@3 & R@5 & R@10 & R@3 & R@5 & R@10 & R@3 & R@5 & R@10 \\
\midrule
Naive Retriever~\cite{historiqa-pellet}              & 35.9 & 44.0 & 54.4 & 14.3 & 19.7 & 28.8 & 47.8 & 57.3 & 68.5 \\
BM25 (lexical)~\cite{bm25}                    & 23.1 & 29.3 & 37.8 &  4.6 &  8.3 & 13.1 & 33.3 & 40.9 & 51.4 \\
HippoRAGv2~\cite{hipporagv2}                  & 31.7 & 41.9 & 53.0 & 11.6 & 17.2 & 24.7 & 42.7 & 55.4 & 68.6 \\
LinearRAG~\cite{linearrag_2025}               & 18.2 & 27.0 & 40.1 & 6.2 & 10.4 & 20.5 & 25.1 & 36.5 & 50.9 \\
\midrule
S-RAG: QRe                                    & 48.0 & 57.0 & 68.4 & \textbf{48.4} & \textbf{56.7} & \underline{68.3} & 47.8 & 57.2 & 68.5 \\
S-RAG: UMS-Retriever                          & \underline{50.2} & \underline{57.6} & \underline{69.4} & 40.8 & 53.3 & 60.7 & \textbf{55.4} & \textbf{60.0} & \textbf{74.2} \\
S-RAG: QRe \& UMS-Retriever                   & \textbf{52.3} & \textbf{58.4} & \textbf{72.3} & \underline{46.7} & \underline{55.6} & \textbf{69.3} & \textbf{55.4} & \textbf{60.0} & \underline{73.9} \\
\bottomrule
\end{tabular*}
\caption{Recall@$k$ on multi-hop questions, with breakdown by question type.
\textbf{S-RAG}: this work. \textbf{QRe}: supervised source pre-selection.
\textbf{UMS}: equal-quota multi-source retrieval.
\textbf{BM25}: lexical baseline using the same chunking as the Naive Retriever.
\textbf{Bold} = best; \underline{underline} = second-best per column (Overall).}
\label{tab:retrieval-results}
\end{table*}

\paragraph{Overall (RQ1, RQ2).}
All three RAG configurations beat the naive, sparse retrieval and graph-based RAG baselines at every cut-off. The combined QRe\,\&\,UMS reaches Recall@3 of 52.3
($+16.4$~\% points over Naive, $+20.6\%$ over HippoRAGv2), with the gap still being high at $k{=}10$.

\paragraph{Performances by question type.}
The two strategies fix different failures. On cross-newspaper
questions, the Naive Retriever collapses (14.3 @3) as \textit{Les Débats}
floods the top-$k$ with parliamentary chunks unrelated to the query. QRe
removes this contamination and more than triples Recall@3 (48.4), while UMS alone
is less effective (40.8) because it still reserves a third of the budget
for \textit{Les Débats}. On \emph{newspaper $\leftrightarrow$ Débats} questions
the asymmetry reverses: the Naive Retriever already performs well (47.8),
QRe is neutral (the routing map keeps all three sources active), and UMS
delivers the main lift (55.4) by guaranteeing parliamentary representation
that would otherwise be crowded out.

\paragraph{Complementarity (RQ2).}
Neither routing alone nor uniform allocation alone accounts for the full
gain: each component contributes independently, and their combination is
strictly dominant overall on every metric. QRe and UMS target structurally distinct
failure modes : source contamination versus source starvation that can co-occur on a single question, which is why combining them
dominates either component in isolation.


\paragraph{Comparison to other methods.}
We include BM25~\cite{bm25} as a sparse-retrieval baseline using the same chunking
as the Naive Retriever. BM25 underperforms the dense baselines on every
cut-off and on both question types, with gap on cross-newspaper questions of 4.6 vs.\ 14.3 Recall@3 against the
Naive Retriever. Lexical matching is particularly ill-suited to this
benchmark: questions are paraphrased rather than quoted from the
articles, and the OCR'd 19th-century newspaper text introduces surface
variation (orthographic and typographic) that dense embeddings absorb
but term-frequency matching does not. We therefore retain dense
retrieval as the underlying index for all Structured retriever variants.

Despite its graph-based index, HippoRAG v2 trails the Naive Retriever on cross-newspaper questions (11.6 vs.\ 14.3 @3) and underperforms all structured retriever variants across every metric. LinearRAG~\cite{linearrag_2025} is the weakest baseline
overall (Recall@3 of 18.2 vs.\ 35.9 for the Naive Retriever) and degrades
sharpest on cross-newspaper questions (6.2 @3, 10.4 @5, 20.5 @10). Even on the
more favourable \emph{newspaper $\leftrightarrow$ D\'ebats} split it remains
below the Naive Retriever at every cut-off (25.1 / 36.5 / 50.9 vs.\ 47.8 / 57.3 / 68.5). Those results suggests that graphs are poorly constructed and that propagation breaks down when evidence is spread across heterogeneous corpora marked by OCR errors and inconsistent chunk sizes.

\subsection{Task/Query-Conditioned Chunking}
\label{sec:indexing-results}

We now analyse whether per-cluster selection of a chunking strategy improves
retrieval over a single global best. The comparison of the 14 candidate
strategies under the four retrievers is in Appendix~\ref{subsec:appendix-global-strategy-comparison}
(Table~\ref{tab:recall-coverage-k3}).
The realised gain on the held-out test set, decomposed by cluster and by
assigned strategy, is reported in
Tables~\ref{tab:chunking-wins-test-cluster} and~\ref{tab:chunking-wins-summary}
(Appendix~\ref{subsec:appendix-cluster-wins}), strategy-level qualitative wins are discussed in
Appendix~\ref{sec:strategy-wins}. A non-trivial subset of clusters
benefits from a strictly different strategy than the global baseline,
motivating the held-out transfer evaluation below.

\paragraph{Transfer to unseen queries (RQ3).}
% The oracle analysis shows the potential for strong retrieval gains when
% the appropriate chunking strategy is chosen, but it assumes the best
% strategy is known in advance. We test the realised gain using the
% inference procedure described in Section~\ref{sec:indexing-inference}.
Results are reported in Table~\ref{tab:tc-inference-merged} (rows
\textit{+\,TC-Chunking}) alongside the metadata-axis variant
discussed below. Under the Naive Retriever the inference pipeline realises
$+5.4$~\% points Coverage@3 and $+6.1$~\% points Recall@5 ($p<0.05$), confirming
that a single low-cost classification step at query time delivers
meaningful gain. Gains are negative but non significant when Qre is active (at most $-0.4$ points) and are negligible and non-significant under UMS retriever (at most $+0.2$~\% points, $p>0.05$ for all routing
configurations), indicating that task-conditioned chunking and source
routing address entirely overlapping failure modes: the gains chunking
provides for the naive retriever are fully recovered by routing alone.

% \paragraph{Mechanism: where TC-Chunking actually helps.}
% A post-hoc decomposition of the Naive-Retriever gain by gold location
% clarifies why TC-Chunking and QRe overlap. Splitting the 875 test
% questions by whether at least one gold passage lies in \textit{Les
% D{\'e}bats} (the long-form parliamentary source, $n{=}561$, 64\%) or
% only in the two newspapers ($n{=}314$, 36\%), the $+5.4$~pp average
% Coverage@3 gain decomposes into $-0.5$~pp on gold-in-D{\'e}bats
% questions and $+16.2$~pp on the remaining 36\%. The benefit of
% TC-Chunking is therefore not better \textit{D{\'e}bats} retrieval but
% \textit{less aggressive competition} from D{\'e}bats chunks in the
% top-$k$, freeing slots for the true newspaper gold. QRe achieves the
% same outcome more directly by filtering D{\'e}bats out of the candidate
% pool whenever the classifier predicts a newspaper-only question; once
% this pre-filter is in place, the noise-suppression channel that
% TC-Chunking exploits is no longer available, and its incremental Δ
% collapses to zero. We confirm this is structural rather than a tuning
% artefact: restricting the QRe evaluation to the subset where
% TC-Chunking can still mechanically act (gold in \textit{Les D{\'e}bats}
% \textit{and} \textit{Les D{\'e}bats} retained by the router; 64\% of
% test, $n{=}283\pm3$ per seed) yields Δ Coverage@3 of $-0.002\pm0.002$,
% and increasing the strategy-selection score-$k$ from 3 to 9 leaves the
% QRe Δ within $\pm 0.001$. Table~\ref{tab:tc-chunking-by-gold-location}
% summarises the decomposition.

% \begin{table}[htbp]
% \centering
% \small
% \setlength{\tabcolsep}{5pt}
% \renewcommand{\arraystretch}{1.15}
% \begin{tabular}{l r ccc}
% \toprule
% \textbf{Subset} & $n$
%   & $\Delta$\,@3 & $\Delta$\,@5 & $\Delta$\,@10 \\
% \midrule
% Full test set                & 875        & $+5.4$  & $+6.3$  & $+6.2$  \\
% \;\;Gold in \textit{Les Débats}     & 561 (64\%) & $-0.5$  & $-0.8$  & $-0.8$  \\
% \;\;Gold not in \textit{Les Débats} & 314 (36\%) & $\mathbf{+16.2}$ & $\mathbf{+19.3}$ & $\mathbf{+19.1}$ \\
% \bottomrule
% \end{tabular}
% \caption{Naive-Retriever Δ Coverage@$k$ (pp, TC-Chunking $-$ Global
% Baseline) on the held-out test set, decomposed by whether at least one
% gold passage lies in \textit{Les Débats}. The aggregate $+5.4$~pp gain
% is carried entirely by the 36\% of questions whose gold lives in the
% two newspapers; on D{\'e}bats-gold questions TC-Chunking is mildly
% detrimental. QRe filters \textit{Les D{\'e}bats} for these same 36\%,
% absorbing the gain and reducing the QRe Δ to zero.}
% \label{tab:tc-chunking-by-gold-location}
% \end{table}

\paragraph{Gain vanishing under routing.}
Decomposing the $+5.4$~\% points Naive-Retriever Coverage@3 gain by gold
location shows it is carried entirely by the questions whose
gold lies only in the newspapers ($+16.2$~\% points), while questions with
gold in \textit{Les D\'ebats} see $-0.5$~\% points. TC-Chunking therefore helps
not by improving \textit{D\'ebats} retrieval but by making \textit{D\'ebats}
chunks less competitive in the top-$k$, freeing slots for the true
newspaper gold. QRe achieves the same outcome more directly filtering
\textit{D\'ebats} out of the candidate pool whenever the classifier
predicts a newspaper-only question which is why the TC-Chunking
$\Delta$ collapses to $\leq 0.2$~\% points once routing is active. Restricted
to the subset where TC-Chunking can still mechanically act
(gold in \textit{Les D\'ebats} \emph{and} \textit{Les D\'ebats} kept by
the router), $\Delta$ Coverage@3 is $-0.002\pm0.002$, and varying the
strategy-selection score-$k$ from 3 to 9 changes it by at most $\pm 0.001$.
Full decomposition and a residual-gain analysis are in
Appendix~\ref{sec:appendix-tc-decomposition}.


% \begin{table*}[htbp]
% \centering
% \small
% \setlength{\tabcolsep}{4pt}
% \begin{tabular}{ll cc cc}
% \toprule
% & & \multicolumn{2}{c}{\textbf{Coverage}} & \multicolumn{2}{c}{\textbf{Recall}} \\
% \cmidrule(lr){3-4} \cmidrule(lr){5-6}
% \textbf{Retrieval} & \textbf{Chunking} & \textbf{@3} & \textbf{@5} & \textbf{@3} & \textbf{@5} \\
% \midrule
% \multirow{3}{*}{Naive}
%   & Global Baseline  & 0.390 & 0.457 & 0.390 & 0.457 \\
%   & Task-Conditioned & \textbf{0.465} & \textbf{0.538} & \textbf{0.466} & \textbf{0.537} \\
%   & $\Delta$         & \cellcolor{db1}{+0.075} & \cellcolor{db1}{+0.081} & \cellcolor{db1}{+0.076} & \cellcolor{db1}{+0.080} \\
% \midrule
% \multirow{3}{*}{QRe}
%   & Global Baseline  & 0.495 & 0.578 & 0.500 & \textbf{0.580} \\
%   & Task-Conditioned & \textbf{0.511} & \textbf{0.581} & \textbf{0.511} & 0.575 \\
%   & $\Delta$         & \cellcolor{db1}{+0.016} & +0.003 & \cellcolor{db1}{+0.011} & $-$0.005 \\
% \midrule
% \multirow{3}{*}{UMS}
%   & Global Baseline  & 0.534 & 0.598 & \textbf{0.534} & 0.600 \\
%   & Task-Conditioned & \textbf{0.543} & \textbf{0.610} & \textbf{0.534} & \textbf{0.603} \\
%   & $\Delta$         & \cellcolor{db1}{+0.009} & \cellcolor{db1}{+0.012} & +0.000 & \cellcolor{db1}{+0.003} \\
% \midrule
% \multirow{3}{*}{QRe \& UMS}
%   & Global Baseline  & 0.546 & 0.607 & \textbf{0.548} & 0.610 \\
%   & Task-Conditioned & \textbf{0.555} & \textbf{0.619} & \textbf{0.548} & \textbf{0.612} \\
%   & $\Delta$         & \cellcolor{db1}{+0.009} & \cellcolor{db1}{+0.012} & +0.000 & \cellcolor{db1}{+0.002} \\
% \bottomrule
% \end{tabular}
% \caption{Coverage@$k$ and Recall@$k$ on the held-out test set ($n{=}221$):
% Global Baseline (\texttt{10000\_hierarchical}) vs.\ Task-Conditioned chunking.
% Task-conditioned gains are largest under the Naive Retriever and shrink
% substantially once routing is active.}
% \label{tab:tc-inference}
% \end{table*}


\subsection{Task-Conditioned Metadata Indexing}
\label{sec:metadata-results}
% Applying the same protocol to the orthogonal metadata axis,
% %(Section~\ref{sec:metadata-methodology})
% we hold chunking fixed at
% \texttt{10000\_hierarchical} and let each cluster pick its best metadata
% strategy on \textit{Les Débats}.
The realised gain on the held-out test set, decomposed by cluster and by
assigned strategy, is reported in
Tables~\ref{tab:metadata-wins-test-cluster} and~\ref{tab:metadata-wins-test-strategy}
(Appendix~\ref{subsec:appendix-metadata-wins}).

\paragraph{Transfer to unseen queries (RQ4).}

Table~\ref{tab:tc-inference-merged} reports
Coverage@$k$ and Recall@$k$ across the four retrievers (rows
\textit{+\,TC-Metadata}). Unlike chunking, the metadata gain persists under UMS: the gain is
largest under the Naive Retriever ($+8.9$~\% points Coverage@3, $+9.0$~\% points
Recall@3; $p<0.05$), remains substantial under UMS ($+6.0$~\% points Cov@3,
$+6.0$~\% points Rec@3; $p<0.05$) where parliamentary chunks occupy a
guaranteed retrieval quota and filtering noisy document types directly
improves that quota, vanishes under QRe ($0.0$~\% points) which already
pre-selects away from noisy sources, and is absent under
QRe\,\&\,UMS ($0.0$~\% points).
%On the held-out set only three distinct
%strategies are ever assigned in practice (\texttt{baseline}: 15 clusters,
%\texttt{exclude\_noisy}: 5, \texttt{rerank\_type\_boost}: 1).

We also evaluated the joint $14{\times}14$ chunking$\times$metadata
strategy space; the joint variant matches but never strictly improves on
TC-Metadata under any retriever, so we defer its methodology, results,
and non-significance analysis to
Appendix~\ref{sec:appendix-joint}.

% \begin{table*}[htbp]
% \centering
% \small
% \setlength{\tabcolsep}{4pt}
% \begin{tabular}{ll cc cc}
% \toprule
% & & \multicolumn{2}{c}{\textbf{Coverage}} & \multicolumn{2}{c}{\textbf{Recall}} \\
% \cmidrule(lr){3-4} \cmidrule(lr){5-6}
% \textbf{Retrieval} & \textbf{Metadata} & \textbf{@3} & \textbf{@5} & \textbf{@3} & \textbf{@5} \\
% \midrule
% \multirow{3}{*}{Naive}
%   & Global Baseline  & 0.390 & 0.457 & 0.390 & 0.457 \\
%   & Task-Conditioned & \textbf{0.498} & \textbf{0.573} & \textbf{0.502} & \textbf{0.575} \\
%   & $\Delta$         & \cellcolor{db1}{+0.108} & \cellcolor{db1}{+0.116} & \cellcolor{db1}{+0.112} & \cellcolor{db1}{+0.118} \\
% \midrule
% \multirow{3}{*}{QRe}
%   & Global Baseline  & 0.495 & 0.578 & 0.500 & 0.580 \\
%   & Task-Conditioned & \textbf{0.505} & \textbf{0.582} & \textbf{0.509} & \textbf{0.584} \\
%   & $\Delta$         & \cellcolor{db1}{+0.010} & \cellcolor{db1}{+0.004} & \cellcolor{db1}{+0.009} & \cellcolor{db1}{+0.004} \\
% \midrule
% \multirow{3}{*}{UMS}
%   & Global Baseline  & 0.534 & 0.598 & 0.534 & 0.600 \\
%   & Task-Conditioned & \textbf{0.559} & \textbf{0.607} & \textbf{0.562} & \textbf{0.610} \\
%   & $\Delta$         & \cellcolor{db1}{+0.025} & \cellcolor{db1}{+0.009} & \cellcolor{db1}{+0.028} & \cellcolor{db1}{+0.010} \\
% \midrule
% \multirow{3}{*}{QRe \& UMS}
%   & Global Baseline  & \textbf{0.546} & \textbf{0.607} & \textbf{0.548} & \textbf{0.610} \\
%   & Task-Conditioned & 0.546 & 0.607 & 0.548 & 0.610 \\
%   & $\Delta$         & +0.000 & +0.000 & +0.000 & +0.000 \\
% \bottomrule
% \end{tabular}
% \caption{Coverage@$k$ and Recall@$k$ on the held-out test set
% ($n_{\text{test}}{=}219$): Global Baseline vs.\ Task-Conditioned metadata
% indexing on \textit{Les Débats}. $k_0{=}6$ for UMS-based routings.}
% \label{tab:tc-meta-inference}
% \end{table*}

\begin{table}[htbp]
\centering
\small
\setlength{\tabcolsep}{4pt}
\renewcommand{\arraystretch}{1.1}
\resizebox{\columnwidth}{!}{%
\begin{tabular}{ll cc cc}
\toprule
\textbf{Retriever} & \textbf{Indexing strategy}
  & \multicolumn{2}{c}{\textbf{Coverage (\%)}}
  & \multicolumn{2}{c}{\textbf{Recall (\%)}} \\
\cmidrule(lr){3-4} \cmidrule(lr){5-6}
 & & @3 & @5 & @3 & @5 \\
\midrule
\multirow{3}{*}{Naive}
  & Global Baseline                 & 36.9               & 45.4               & 36.9               & 45.6               \\
  & \;+\,TC-Chunking                & 42.3$^{*}$         & 51.7$^{*}$         & 42.3$^{*}$         & 51.7$^{*}$         \\
  & \;+\,TC-Metadata                & \best{45.8}$^{*}$  & \best{55.2}$^{*}$  & \best{45.9}$^{*}$  & \best{55.3}$^{*}$  \\
\midrule
\multirow{3}{*}{QRe}
  & Global Baseline                 & 48.5               & 58.2               & 48.7               & 58.4               \\
  & \;+\,TC-Chunking                & \uline{48.5}$^{\text{ns}}$ & \uline{57.9}$^{\text{ns}}$ & \uline{48.5}$^{\text{ns}}$ & \uline{58.0}$^{\text{ns}}$ \\
  & \;+\,TC-Metadata                & \best{48.5}$^{\text{ns}}$ & \best{58.2}$^{\text{ns}}$ & \best{48.7}$^{\text{ns}}$ & \best{58.4}$^{\text{ns}}$ \\
\midrule
\multirow{3}{*}{UMS}
  & Global Baseline                 & 46.8               & 58.7               & 47.0               & 58.8               \\
  & \;+\,TC-Chunking                & 47.0$^{\text{ns}}$ & 58.5$^{\text{ns}}$ & 47.2$^{\text{ns}}$ & 58.7$^{\text{ns}}$ \\
  & \;+\,TC-Metadata                & \best{52.8}$^{*}$  & \best{63.2}$^{*}$  & \best{53.0}$^{*}$  & \best{63.4}$^{*}$  \\
\midrule
\multirow{3}{*}{QRe \& UMS}
  & Global Baseline                 & 54.7               & 64.5               & 54.9               & 64.7               \\
  & \;+\,TC-Chunking                & \uline{54.7}$^{\text{ns}}$ & \uline{64.5}$^{\text{ns}}$ & \uline{54.9}$^{\text{ns}}$ & \uline{64.7}$^{\text{ns}}$ \\
  & \;+\,TC-Metadata                & \best{54.7}$^{\text{ns}}$ & \best{64.5}$^{\text{ns}}$ & \best{54.9}$^{\text{ns}}$ & \best{64.7}$^{\text{ns}}$ \\
\bottomrule
\end{tabular}
}
\caption{Coverage@$k$ and Recall@$k$ on the held-out test set
($n{=}438$, 5 seeds), comparing each retriever's Global Baseline
(\texttt{10000\_hier.} chunking, no metadata filter, applied uniformly
to every query) to two task-conditioned variants. Under
\textit{+\,TC-Chunking}, every test question is first assigned to its
nearest question cluster, and the chunking strategy used to retrieve
its chunks is the one found optimal for that cluster on the training
split (selected from the 14 candidates of
Appendix~\ref{sec:appendix-chunking}); chunking is therefore
per-cluster, i.e.\ effectively per-question. \textit{+\,TC-Metadata}
applies the same cluster-routing logic to the metadata axis instead,
holding chunking fixed at the baseline and letting each cluster pick
its optimal metadata filter / reranker from the
14 candidates. \textbf{Bold} / \uline{underline} = best / second within block.
$^{*}p<0.05$ (Bonferroni paired $t$-test); ns = not significant.}
\label{tab:tc-inference-merged}
\end{table} in main paper
% ============================================================

\definecolor{sectiongray}{gray}{0.95}

\section{Results}
\label{sec:results}

We address four questions:
\textbf{RQ1} : does query-conditioned source selection improve retrieval over a single-index baseline?
\textbf{RQ2} : are query routing (QRe) and uniform multi-source allocation (UMS) complementary remedies for size-imbalanced retrieval?
\textbf{RQ3} : does query-conditioned chunking optimisation improve retrieval over a global baseline?
\textbf{RQ4} : does query-conditioned \emph{metadata indexing} (filters and rerankers on \textit{Les Débats}) improve retrieval over a global baseline?

\subsection{Retrieval Strategies}
\label{sec:retrieval-results}

We evaluate the four retrieval strategies of
Section~\ref{sec:retrieval-strategies} on the multi-hop subset of
HistoriQA-ThirdRepublic~\cite{historiqa-pellet}, broken down by question type:
\textit{cross-newspaper} (314 q., both gold documents from newspapers) and
\textit{newspaper $\rightarrow$ Débats} (571 q., one newspaper plus one
parliamentary document). We include comparison to BM25~\cite{bm25} for sparse retrieval and 2 graph-based external baseline: HippoRAGv2~\cite{hipporagv2} and LinearRAG~\cite{linearrag_2025}.


\begin{table*}[ht]
\centering
\small
\setlength{\tabcolsep}{4pt}
\begin{tabular*}{\textwidth}{@{\extracolsep{\fill}} l ccc ccc ccc @{}}
\toprule
 & \multicolumn{3}{c}{\textbf{Overall}} & \multicolumn{3}{c}{\textbf{Cross-newspaper} (314 q.)} & \multicolumn{3}{c}{\textbf{Newspaper-Débats} (571 q.)} \\
\cmidrule(lr){2-4} \cmidrule(lr){5-7} \cmidrule(lr){8-10}
\textbf{Configuration} & R@3 & R@5 & R@10 & R@3 & R@5 & R@10 & R@3 & R@5 & R@10 \\
\midrule
Naive Retriever~\cite{historiqa-pellet}              & 35.9 & 44.0 & 54.4 & 14.3 & 19.7 & 28.8 & 47.8 & 57.3 & 68.5 \\
BM25 (lexical)~\cite{bm25}                    & 23.1 & 29.3 & 37.8 &  4.6 &  8.3 & 13.1 & 33.3 & 40.9 & 51.4 \\
HippoRAGv2~\cite{hipporagv2}                  & 31.7 & 41.9 & 53.0 & 11.6 & 17.2 & 24.7 & 42.7 & 55.4 & 68.6 \\
LinearRAG~\cite{linearrag_2025}               & 18.2 & 27.0 & 40.1 & 6.2 & 10.4 & 20.5 & 25.1 & 36.5 & 50.9 \\
\midrule
S-RAG: QRe                                    & 48.0 & 57.0 & 68.4 & \textbf{48.4} & \textbf{56.7} & \underline{68.3} & 47.8 & 57.2 & 68.5 \\
S-RAG: UMS-Retriever                          & \underline{50.2} & \underline{57.6} & \underline{69.4} & 40.8 & 53.3 & 60.7 & \textbf{55.4} & \textbf{60.0} & \textbf{74.2} \\
S-RAG: QRe \& UMS-Retriever                   & \textbf{52.3} & \textbf{58.4} & \textbf{72.3} & \underline{46.7} & \underline{55.6} & \textbf{69.3} & \textbf{55.4} & \textbf{60.0} & \underline{73.9} \\
\bottomrule
\end{tabular*}
\caption{Recall@$k$ on multi-hop questions, with breakdown by question type.
\textbf{S-RAG}: this work. \textbf{QRe}: supervised source pre-selection.
\textbf{UMS}: equal-quota multi-source retrieval.
\textbf{BM25}: lexical baseline using the same chunking as the Naive Retriever.
\textbf{Bold} = best; \underline{underline} = second-best per column (Overall).}
\label{tab:retrieval-results}
\end{table*}

\paragraph{Overall (RQ1, RQ2).}
All three RAG configurations beat the naive, sparse retrieval and graph-based RAG baselines at every cut-off. The combined QRe\,\&\,UMS reaches Recall@3 of 52.3
($+16.4$~\% points over Naive, $+20.6\%$ over HippoRAGv2), with the gap still being high at $k{=}10$.

\paragraph{Performances by question type.}
The two strategies fix different failures. On cross-newspaper
questions, the Naive Retriever collapses (14.3 @3) as \textit{Les Débats}
floods the top-$k$ with parliamentary chunks unrelated to the query. QRe
removes this contamination and more than triples Recall@3 (48.4), while UMS alone
is less effective (40.8) because it still reserves a third of the budget
for \textit{Les Débats}. On \emph{newspaper $\leftrightarrow$ Débats} questions
the asymmetry reverses: the Naive Retriever already performs well (47.8),
QRe is neutral (the routing map keeps all three sources active), and UMS
delivers the main lift (55.4) by guaranteeing parliamentary representation
that would otherwise be crowded out.

\paragraph{Complementarity (RQ2).}
Neither routing alone nor uniform allocation alone accounts for the full
gain: each component contributes independently, and their combination is
strictly dominant overall on every metric. QRe and UMS target structurally distinct
failure modes : source contamination versus source starvation that can co-occur on a single question, which is why combining them
dominates either component in isolation.


\paragraph{Comparison to other methods.}
We include BM25~\cite{bm25} as a sparse-retrieval baseline using the same chunking
as the Naive Retriever. BM25 underperforms the dense baselines on every
cut-off and on both question types, with gap on cross-newspaper questions of 4.6 vs.\ 14.3 Recall@3 against the
Naive Retriever. Lexical matching is particularly ill-suited to this
benchmark: questions are paraphrased rather than quoted from the
articles, and the OCR'd 19th-century newspaper text introduces surface
variation (orthographic and typographic) that dense embeddings absorb
but term-frequency matching does not. We therefore retain dense
retrieval as the underlying index for all Structured retriever variants.

Despite its graph-based index, HippoRAG v2 trails the Naive Retriever on cross-newspaper questions (11.6 vs.\ 14.3 @3) and underperforms all structured retriever variants across every metric. LinearRAG~\cite{linearrag_2025} is the weakest baseline
overall (Recall@3 of 18.2 vs.\ 35.9 for the Naive Retriever) and degrades
sharpest on cross-newspaper questions (6.2 @3, 10.4 @5, 20.5 @10). Even on the
more favourable \emph{newspaper $\leftrightarrow$ D\'ebats} split it remains
below the Naive Retriever at every cut-off (25.1 / 36.5 / 50.9 vs.\ 47.8 / 57.3 / 68.5). Those results suggests that graphs are poorly constructed and that propagation breaks down when evidence is spread across heterogeneous corpora marked by OCR errors and inconsistent chunk sizes.

\subsection{Task/Query-Conditioned Chunking}
\label{sec:indexing-results}

We now analyse whether per-cluster selection of a chunking strategy improves
retrieval over a single global best. The comparison of the 14 candidate
strategies under the four retrievers is in Appendix~\ref{subsec:appendix-global-strategy-comparison}
(Table~\ref{tab:recall-coverage-k3}).
The realised gain on the held-out test set, decomposed by cluster and by
assigned strategy, is reported in
Tables~\ref{tab:chunking-wins-test-cluster} and~\ref{tab:chunking-wins-summary}
(Appendix~\ref{subsec:appendix-cluster-wins}), strategy-level qualitative wins are discussed in
Appendix~\ref{sec:strategy-wins}. A non-trivial subset of clusters
benefits from a strictly different strategy than the global baseline,
motivating the held-out transfer evaluation below.

\paragraph{Transfer to unseen queries (RQ3).}
% The oracle analysis shows the potential for strong retrieval gains when
% the appropriate chunking strategy is chosen, but it assumes the best
% strategy is known in advance. We test the realised gain using the
% inference procedure described in Section~\ref{sec:indexing-inference}.
Results are reported in Table~\ref{tab:tc-inference-merged} (rows
\textit{+\,TC-Chunking}) alongside the metadata-axis variant
discussed below. Under the Naive Retriever the inference pipeline realises
$+5.4$~\% points Coverage@3 and $+6.1$~\% points Recall@5 ($p<0.05$), confirming
that a single low-cost classification step at query time delivers
meaningful gain. Gains are negative but non significant when Qre is active (at most $-0.4$ points) and are negligible and non-significant under UMS retriever (at most $+0.2$~\% points, $p>0.05$ for all routing
configurations), indicating that task-conditioned chunking and source
routing address entirely overlapping failure modes: the gains chunking
provides for the naive retriever are fully recovered by routing alone.

% \paragraph{Mechanism: where TC-Chunking actually helps.}
% A post-hoc decomposition of the Naive-Retriever gain by gold location
% clarifies why TC-Chunking and QRe overlap. Splitting the 875 test
% questions by whether at least one gold passage lies in \textit{Les
% D{\'e}bats} (the long-form parliamentary source, $n{=}561$, 64\%) or
% only in the two newspapers ($n{=}314$, 36\%), the $+5.4$~pp average
% Coverage@3 gain decomposes into $-0.5$~pp on gold-in-D{\'e}bats
% questions and $+16.2$~pp on the remaining 36\%. The benefit of
% TC-Chunking is therefore not better \textit{D{\'e}bats} retrieval but
% \textit{less aggressive competition} from D{\'e}bats chunks in the
% top-$k$, freeing slots for the true newspaper gold. QRe achieves the
% same outcome more directly by filtering D{\'e}bats out of the candidate
% pool whenever the classifier predicts a newspaper-only question; once
% this pre-filter is in place, the noise-suppression channel that
% TC-Chunking exploits is no longer available, and its incremental Δ
% collapses to zero. We confirm this is structural rather than a tuning
% artefact: restricting the QRe evaluation to the subset where
% TC-Chunking can still mechanically act (gold in \textit{Les D{\'e}bats}
% \textit{and} \textit{Les D{\'e}bats} retained by the router; 64\% of
% test, $n{=}283\pm3$ per seed) yields Δ Coverage@3 of $-0.002\pm0.002$,
% and increasing the strategy-selection score-$k$ from 3 to 9 leaves the
% QRe Δ within $\pm 0.001$. Table~\ref{tab:tc-chunking-by-gold-location}
% summarises the decomposition.

% \begin{table}[htbp]
% \centering
% \small
% \setlength{\tabcolsep}{5pt}
% \renewcommand{\arraystretch}{1.15}
% \begin{tabular}{l r ccc}
% \toprule
% \textbf{Subset} & $n$
%   & $\Delta$\,@3 & $\Delta$\,@5 & $\Delta$\,@10 \\
% \midrule
% Full test set                & 875        & $+5.4$  & $+6.3$  & $+6.2$  \\
% \;\;Gold in \textit{Les Débats}     & 561 (64\%) & $-0.5$  & $-0.8$  & $-0.8$  \\
% \;\;Gold not in \textit{Les Débats} & 314 (36\%) & $\mathbf{+16.2}$ & $\mathbf{+19.3}$ & $\mathbf{+19.1}$ \\
% \bottomrule
% \end{tabular}
% \caption{Naive-Retriever Δ Coverage@$k$ (pp, TC-Chunking $-$ Global
% Baseline) on the held-out test set, decomposed by whether at least one
% gold passage lies in \textit{Les Débats}. The aggregate $+5.4$~pp gain
% is carried entirely by the 36\% of questions whose gold lives in the
% two newspapers; on D{\'e}bats-gold questions TC-Chunking is mildly
% detrimental. QRe filters \textit{Les D{\'e}bats} for these same 36\%,
% absorbing the gain and reducing the QRe Δ to zero.}
% \label{tab:tc-chunking-by-gold-location}
% \end{table}

\paragraph{Gain vanishing under routing.}
Decomposing the $+5.4$~\% points Naive-Retriever Coverage@3 gain by gold
location shows it is carried entirely by the questions whose
gold lies only in the newspapers ($+16.2$~\% points), while questions with
gold in \textit{Les D\'ebats} see $-0.5$~\% points. TC-Chunking therefore helps
not by improving \textit{D\'ebats} retrieval but by making \textit{D\'ebats}
chunks less competitive in the top-$k$, freeing slots for the true
newspaper gold. QRe achieves the same outcome more directly filtering
\textit{D\'ebats} out of the candidate pool whenever the classifier
predicts a newspaper-only question which is why the TC-Chunking
$\Delta$ collapses to $\leq 0.2$~\% points once routing is active. Restricted
to the subset where TC-Chunking can still mechanically act
(gold in \textit{Les D\'ebats} \emph{and} \textit{Les D\'ebats} kept by
the router), $\Delta$ Coverage@3 is $-0.002\pm0.002$, and varying the
strategy-selection score-$k$ from 3 to 9 changes it by at most $\pm 0.001$.
Full decomposition and a residual-gain analysis are in
Appendix~\ref{sec:appendix-tc-decomposition}.


% \begin{table*}[htbp]
% \centering
% \small
% \setlength{\tabcolsep}{4pt}
% \begin{tabular}{ll cc cc}
% \toprule
% & & \multicolumn{2}{c}{\textbf{Coverage}} & \multicolumn{2}{c}{\textbf{Recall}} \\
% \cmidrule(lr){3-4} \cmidrule(lr){5-6}
% \textbf{Retrieval} & \textbf{Chunking} & \textbf{@3} & \textbf{@5} & \textbf{@3} & \textbf{@5} \\
% \midrule
% \multirow{3}{*}{Naive}
%   & Global Baseline  & 0.390 & 0.457 & 0.390 & 0.457 \\
%   & Task-Conditioned & \textbf{0.465} & \textbf{0.538} & \textbf{0.466} & \textbf{0.537} \\
%   & $\Delta$         & \cellcolor{db1}{+0.075} & \cellcolor{db1}{+0.081} & \cellcolor{db1}{+0.076} & \cellcolor{db1}{+0.080} \\
% \midrule
% \multirow{3}{*}{QRe}
%   & Global Baseline  & 0.495 & 0.578 & 0.500 & \textbf{0.580} \\
%   & Task-Conditioned & \textbf{0.511} & \textbf{0.581} & \textbf{0.511} & 0.575 \\
%   & $\Delta$         & \cellcolor{db1}{+0.016} & +0.003 & \cellcolor{db1}{+0.011} & $-$0.005 \\
% \midrule
% \multirow{3}{*}{UMS}
%   & Global Baseline  & 0.534 & 0.598 & \textbf{0.534} & 0.600 \\
%   & Task-Conditioned & \textbf{0.543} & \textbf{0.610} & \textbf{0.534} & \textbf{0.603} \\
%   & $\Delta$         & \cellcolor{db1}{+0.009} & \cellcolor{db1}{+0.012} & +0.000 & \cellcolor{db1}{+0.003} \\
% \midrule
% \multirow{3}{*}{QRe \& UMS}
%   & Global Baseline  & 0.546 & 0.607 & \textbf{0.548} & 0.610 \\
%   & Task-Conditioned & \textbf{0.555} & \textbf{0.619} & \textbf{0.548} & \textbf{0.612} \\
%   & $\Delta$         & \cellcolor{db1}{+0.009} & \cellcolor{db1}{+0.012} & +0.000 & \cellcolor{db1}{+0.002} \\
% \bottomrule
% \end{tabular}
% \caption{Coverage@$k$ and Recall@$k$ on the held-out test set ($n{=}221$):
% Global Baseline (\texttt{10000\_hierarchical}) vs.\ Task-Conditioned chunking.
% Task-conditioned gains are largest under the Naive Retriever and shrink
% substantially once routing is active.}
% \label{tab:tc-inference}
% \end{table*}


\subsection{Task-Conditioned Metadata Indexing}
\label{sec:metadata-results}
% Applying the same protocol to the orthogonal metadata axis,
% %(Section~\ref{sec:metadata-methodology})
% we hold chunking fixed at
% \texttt{10000\_hierarchical} and let each cluster pick its best metadata
% strategy on \textit{Les Débats}.
The realised gain on the held-out test set, decomposed by cluster and by
assigned strategy, is reported in
Tables~\ref{tab:metadata-wins-test-cluster} and~\ref{tab:metadata-wins-test-strategy}
(Appendix~\ref{subsec:appendix-metadata-wins}).

\paragraph{Transfer to unseen queries (RQ4).}

Table~\ref{tab:tc-inference-merged} reports
Coverage@$k$ and Recall@$k$ across the four retrievers (rows
\textit{+\,TC-Metadata}). Unlike chunking, the metadata gain persists under UMS: the gain is
largest under the Naive Retriever ($+8.9$~\% points Coverage@3, $+9.0$~\% points
Recall@3; $p<0.05$), remains substantial under UMS ($+6.0$~\% points Cov@3,
$+6.0$~\% points Rec@3; $p<0.05$) where parliamentary chunks occupy a
guaranteed retrieval quota and filtering noisy document types directly
improves that quota, vanishes under QRe ($0.0$~\% points) which already
pre-selects away from noisy sources, and is absent under
QRe\,\&\,UMS ($0.0$~\% points).
%On the held-out set only three distinct
%strategies are ever assigned in practice (\texttt{baseline}: 15 clusters,
%\texttt{exclude\_noisy}: 5, \texttt{rerank\_type\_boost}: 1).

We also evaluated the joint $14{\times}14$ chunking$\times$metadata
strategy space; the joint variant matches but never strictly improves on
TC-Metadata under any retriever, so we defer its methodology, results,
and non-significance analysis to
Appendix~\ref{sec:appendix-joint}.

% \begin{table*}[htbp]
% \centering
% \small
% \setlength{\tabcolsep}{4pt}
% \begin{tabular}{ll cc cc}
% \toprule
% & & \multicolumn{2}{c}{\textbf{Coverage}} & \multicolumn{2}{c}{\textbf{Recall}} \\
% \cmidrule(lr){3-4} \cmidrule(lr){5-6}
% \textbf{Retrieval} & \textbf{Metadata} & \textbf{@3} & \textbf{@5} & \textbf{@3} & \textbf{@5} \\
% \midrule
% \multirow{3}{*}{Naive}
%   & Global Baseline  & 0.390 & 0.457 & 0.390 & 0.457 \\
%   & Task-Conditioned & \textbf{0.498} & \textbf{0.573} & \textbf{0.502} & \textbf{0.575} \\
%   & $\Delta$         & \cellcolor{db1}{+0.108} & \cellcolor{db1}{+0.116} & \cellcolor{db1}{+0.112} & \cellcolor{db1}{+0.118} \\
% \midrule
% \multirow{3}{*}{QRe}
%   & Global Baseline  & 0.495 & 0.578 & 0.500 & 0.580 \\
%   & Task-Conditioned & \textbf{0.505} & \textbf{0.582} & \textbf{0.509} & \textbf{0.584} \\
%   & $\Delta$         & \cellcolor{db1}{+0.010} & \cellcolor{db1}{+0.004} & \cellcolor{db1}{+0.009} & \cellcolor{db1}{+0.004} \\
% \midrule
% \multirow{3}{*}{UMS}
%   & Global Baseline  & 0.534 & 0.598 & 0.534 & 0.600 \\
%   & Task-Conditioned & \textbf{0.559} & \textbf{0.607} & \textbf{0.562} & \textbf{0.610} \\
%   & $\Delta$         & \cellcolor{db1}{+0.025} & \cellcolor{db1}{+0.009} & \cellcolor{db1}{+0.028} & \cellcolor{db1}{+0.010} \\
% \midrule
% \multirow{3}{*}{QRe \& UMS}
%   & Global Baseline  & \textbf{0.546} & \textbf{0.607} & \textbf{0.548} & \textbf{0.610} \\
%   & Task-Conditioned & 0.546 & 0.607 & 0.548 & 0.610 \\
%   & $\Delta$         & +0.000 & +0.000 & +0.000 & +0.000 \\
% \bottomrule
% \end{tabular}
% \caption{Coverage@$k$ and Recall@$k$ on the held-out test set
% ($n_{\text{test}}{=}219$): Global Baseline vs.\ Task-Conditioned metadata
% indexing on \textit{Les Débats}. $k_0{=}6$ for UMS-based routings.}
% \label{tab:tc-meta-inference}
% \end{table*}

\begin{table}[htbp]
\centering
\small
\setlength{\tabcolsep}{4pt}
\renewcommand{\arraystretch}{1.1}
\resizebox{\columnwidth}{!}{%
\begin{tabular}{ll cc cc}
\toprule
\textbf{Retriever} & \textbf{Indexing strategy}
  & \multicolumn{2}{c}{\textbf{Coverage (\%)}}
  & \multicolumn{2}{c}{\textbf{Recall (\%)}} \\
\cmidrule(lr){3-4} \cmidrule(lr){5-6}
 & & @3 & @5 & @3 & @5 \\
\midrule
\multirow{3}{*}{Naive}
  & Global Baseline                 & 36.9               & 45.4               & 36.9               & 45.6               \\
  & \;+\,TC-Chunking                & 42.3$^{*}$         & 51.7$^{*}$         & 42.3$^{*}$         & 51.7$^{*}$         \\
  & \;+\,TC-Metadata                & \best{45.8}$^{*}$  & \best{55.2}$^{*}$  & \best{45.9}$^{*}$  & \best{55.3}$^{*}$  \\
\midrule
\multirow{3}{*}{QRe}
  & Global Baseline                 & 48.5               & 58.2               & 48.7               & 58.4               \\
  & \;+\,TC-Chunking                & \uline{48.5}$^{\text{ns}}$ & \uline{57.9}$^{\text{ns}}$ & \uline{48.5}$^{\text{ns}}$ & \uline{58.0}$^{\text{ns}}$ \\
  & \;+\,TC-Metadata                & \best{48.5}$^{\text{ns}}$ & \best{58.2}$^{\text{ns}}$ & \best{48.7}$^{\text{ns}}$ & \best{58.4}$^{\text{ns}}$ \\
\midrule
\multirow{3}{*}{UMS}
  & Global Baseline                 & 46.8               & 58.7               & 47.0               & 58.8               \\
  & \;+\,TC-Chunking                & 47.0$^{\text{ns}}$ & 58.5$^{\text{ns}}$ & 47.2$^{\text{ns}}$ & 58.7$^{\text{ns}}$ \\
  & \;+\,TC-Metadata                & \best{52.8}$^{*}$  & \best{63.2}$^{*}$  & \best{53.0}$^{*}$  & \best{63.4}$^{*}$  \\
\midrule
\multirow{3}{*}{QRe \& UMS}
  & Global Baseline                 & 54.7               & 64.5               & 54.9               & 64.7               \\
  & \;+\,TC-Chunking                & \uline{54.7}$^{\text{ns}}$ & \uline{64.5}$^{\text{ns}}$ & \uline{54.9}$^{\text{ns}}$ & \uline{64.7}$^{\text{ns}}$ \\
  & \;+\,TC-Metadata                & \best{54.7}$^{\text{ns}}$ & \best{64.5}$^{\text{ns}}$ & \best{54.9}$^{\text{ns}}$ & \best{64.7}$^{\text{ns}}$ \\
\bottomrule
\end{tabular}
}
\caption{Coverage@$k$ and Recall@$k$ on the held-out test set
($n{=}438$, 5 seeds), comparing each retriever's Global Baseline
(\texttt{10000\_hier.} chunking, no metadata filter, applied uniformly
to every query) to two task-conditioned variants. Under
\textit{+\,TC-Chunking}, every test question is first assigned to its
nearest question cluster, and the chunking strategy used to retrieve
its chunks is the one found optimal for that cluster on the training
split (selected from the 14 candidates of
Appendix~\ref{sec:appendix-chunking}); chunking is therefore
per-cluster, i.e.\ effectively per-question. \textit{+\,TC-Metadata}
applies the same cluster-routing logic to the metadata axis instead,
holding chunking fixed at the baseline and letting each cluster pick
its optimal metadata filter / reranker from the
14 candidates. \textbf{Bold} / \uline{underline} = best / second within block.
$^{*}p<0.05$ (Bonferroni paired $t$-test); ns = not significant.}
\label{tab:tc-inference-merged}
\end{table} in main paper
% ============================================================

\definecolor{sectiongray}{gray}{0.95}

\section{Results}
\label{sec:results}

We address four questions:
\textbf{RQ1} : does query-conditioned source selection improve retrieval over a single-index baseline?
\textbf{RQ2} : are query routing (QRe) and uniform multi-source allocation (UMS) complementary remedies for size-imbalanced retrieval?
\textbf{RQ3} : does query-conditioned chunking optimisation improve retrieval over a global baseline?
\textbf{RQ4} : does query-conditioned \emph{metadata indexing} (filters and rerankers on \textit{Les Débats}) improve retrieval over a global baseline?

\subsection{Retrieval Strategies}
\label{sec:retrieval-results}

We evaluate the four retrieval strategies of
Section~\ref{sec:retrieval-strategies} on the multi-hop subset of
HistoriQA-ThirdRepublic~\cite{historiqa-pellet}, broken down by question type:
\textit{cross-newspaper} (314 q., both gold documents from newspapers) and
\textit{newspaper $\rightarrow$ Débats} (571 q., one newspaper plus one
parliamentary document). We include comparison to BM25~\cite{bm25} for sparse retrieval and 2 graph-based external baseline: HippoRAGv2~\cite{hipporagv2} and LinearRAG~\cite{linearrag_2025}.


\begin{table*}[ht]
\centering
\small
\setlength{\tabcolsep}{4pt}
\begin{tabular*}{\textwidth}{@{\extracolsep{\fill}} l ccc ccc ccc @{}}
\toprule
 & \multicolumn{3}{c}{\textbf{Overall}} & \multicolumn{3}{c}{\textbf{Cross-newspaper} (314 q.)} & \multicolumn{3}{c}{\textbf{Newspaper-Débats} (571 q.)} \\
\cmidrule(lr){2-4} \cmidrule(lr){5-7} \cmidrule(lr){8-10}
\textbf{Configuration} & R@3 & R@5 & R@10 & R@3 & R@5 & R@10 & R@3 & R@5 & R@10 \\
\midrule
Naive Retriever~\cite{historiqa-pellet}              & 35.9 & 44.0 & 54.4 & 14.3 & 19.7 & 28.8 & 47.8 & 57.3 & 68.5 \\
BM25 (lexical)~\cite{bm25}                    & 23.1 & 29.3 & 37.8 &  4.6 &  8.3 & 13.1 & 33.3 & 40.9 & 51.4 \\
HippoRAGv2~\cite{hipporagv2}                  & 31.7 & 41.9 & 53.0 & 11.6 & 17.2 & 24.7 & 42.7 & 55.4 & 68.6 \\
LinearRAG~\cite{linearrag_2025}               & 18.2 & 27.0 & 40.1 & 6.2 & 10.4 & 20.5 & 25.1 & 36.5 & 50.9 \\
\midrule
S-RAG: QRe                                    & 48.0 & 57.0 & 68.4 & \textbf{48.4} & \textbf{56.7} & \underline{68.3} & 47.8 & 57.2 & 68.5 \\
S-RAG: UMS-Retriever                          & \underline{50.2} & \underline{57.6} & \underline{69.4} & 40.8 & 53.3 & 60.7 & \textbf{55.4} & \textbf{60.0} & \textbf{74.2} \\
S-RAG: QRe \& UMS-Retriever                   & \textbf{52.3} & \textbf{58.4} & \textbf{72.3} & \underline{46.7} & \underline{55.6} & \textbf{69.3} & \textbf{55.4} & \textbf{60.0} & \underline{73.9} \\
\bottomrule
\end{tabular*}
\caption{Recall@$k$ on multi-hop questions, with breakdown by question type.
\textbf{S-RAG}: this work. \textbf{QRe}: supervised source pre-selection.
\textbf{UMS}: equal-quota multi-source retrieval.
\textbf{BM25}: lexical baseline using the same chunking as the Naive Retriever.
\textbf{Bold} = best; \underline{underline} = second-best per column (Overall).}
\label{tab:retrieval-results}
\end{table*}

\paragraph{Overall (RQ1, RQ2).}
All three RAG configurations beat the naive, sparse retrieval and graph-based RAG baselines at every cut-off. The combined QRe\,\&\,UMS reaches Recall@3 of 52.3
($+16.4$~\% points over Naive, $+20.6\%$ over HippoRAGv2), with the gap still being high at $k{=}10$.

\paragraph{Performances by question type.}
The two strategies fix different failures. On cross-newspaper
questions, the Naive Retriever collapses (14.3 @3) as \textit{Les Débats}
floods the top-$k$ with parliamentary chunks unrelated to the query. QRe
removes this contamination and more than triples Recall@3 (48.4), while UMS alone
is less effective (40.8) because it still reserves a third of the budget
for \textit{Les Débats}. On \emph{newspaper $\leftrightarrow$ Débats} questions
the asymmetry reverses: the Naive Retriever already performs well (47.8),
QRe is neutral (the routing map keeps all three sources active), and UMS
delivers the main lift (55.4) by guaranteeing parliamentary representation
that would otherwise be crowded out.

\paragraph{Complementarity (RQ2).}
Neither routing alone nor uniform allocation alone accounts for the full
gain: each component contributes independently, and their combination is
strictly dominant overall on every metric. QRe and UMS target structurally distinct
failure modes : source contamination versus source starvation that can co-occur on a single question, which is why combining them
dominates either component in isolation.


\paragraph{Comparison to other methods.}
We include BM25~\cite{bm25} as a sparse-retrieval baseline using the same chunking
as the Naive Retriever. BM25 underperforms the dense baselines on every
cut-off and on both question types, with gap on cross-newspaper questions of 4.6 vs.\ 14.3 Recall@3 against the
Naive Retriever. Lexical matching is particularly ill-suited to this
benchmark: questions are paraphrased rather than quoted from the
articles, and the OCR'd 19th-century newspaper text introduces surface
variation (orthographic and typographic) that dense embeddings absorb
but term-frequency matching does not. We therefore retain dense
retrieval as the underlying index for all Structured retriever variants.

Despite its graph-based index, HippoRAG v2 trails the Naive Retriever on cross-newspaper questions (11.6 vs.\ 14.3 @3) and underperforms all structured retriever variants across every metric. LinearRAG~\cite{linearrag_2025} is the weakest baseline
overall (Recall@3 of 18.2 vs.\ 35.9 for the Naive Retriever) and degrades
sharpest on cross-newspaper questions (6.2 @3, 10.4 @5, 20.5 @10). Even on the
more favourable \emph{newspaper $\leftrightarrow$ D\'ebats} split it remains
below the Naive Retriever at every cut-off (25.1 / 36.5 / 50.9 vs.\ 47.8 / 57.3 / 68.5). Those results suggests that graphs are poorly constructed and that propagation breaks down when evidence is spread across heterogeneous corpora marked by OCR errors and inconsistent chunk sizes.

\subsection{Task/Query-Conditioned Chunking}
\label{sec:indexing-results}

We now analyse whether per-cluster selection of a chunking strategy improves
retrieval over a single global best. The comparison of the 14 candidate
strategies under the four retrievers is in Appendix~\ref{subsec:appendix-global-strategy-comparison}
(Table~\ref{tab:recall-coverage-k3}).
The realised gain on the held-out test set, decomposed by cluster and by
assigned strategy, is reported in
Tables~\ref{tab:chunking-wins-test-cluster} and~\ref{tab:chunking-wins-summary}
(Appendix~\ref{subsec:appendix-cluster-wins}), strategy-level qualitative wins are discussed in
Appendix~\ref{sec:strategy-wins}. A non-trivial subset of clusters
benefits from a strictly different strategy than the global baseline,
motivating the held-out transfer evaluation below.

\paragraph{Transfer to unseen queries (RQ3).}
% The oracle analysis shows the potential for strong retrieval gains when
% the appropriate chunking strategy is chosen, but it assumes the best
% strategy is known in advance. We test the realised gain using the
% inference procedure described in Section~\ref{sec:indexing-inference}.
Results are reported in Table~\ref{tab:tc-inference-merged} (rows
\textit{+\,TC-Chunking}) alongside the metadata-axis variant
discussed below. Under the Naive Retriever the inference pipeline realises
$+5.4$~\% points Coverage@3 and $+6.1$~\% points Recall@5 ($p<0.05$), confirming
that a single low-cost classification step at query time delivers
meaningful gain. Gains are negative but non significant when Qre is active (at most $-0.4$ points) and are negligible and non-significant under UMS retriever (at most $+0.2$~\% points, $p>0.05$ for all routing
configurations), indicating that task-conditioned chunking and source
routing address entirely overlapping failure modes: the gains chunking
provides for the naive retriever are fully recovered by routing alone.

% \paragraph{Mechanism: where TC-Chunking actually helps.}
% A post-hoc decomposition of the Naive-Retriever gain by gold location
% clarifies why TC-Chunking and QRe overlap. Splitting the 875 test
% questions by whether at least one gold passage lies in \textit{Les
% D{\'e}bats} (the long-form parliamentary source, $n{=}561$, 64\%) or
% only in the two newspapers ($n{=}314$, 36\%), the $+5.4$~pp average
% Coverage@3 gain decomposes into $-0.5$~pp on gold-in-D{\'e}bats
% questions and $+16.2$~pp on the remaining 36\%. The benefit of
% TC-Chunking is therefore not better \textit{D{\'e}bats} retrieval but
% \textit{less aggressive competition} from D{\'e}bats chunks in the
% top-$k$, freeing slots for the true newspaper gold. QRe achieves the
% same outcome more directly by filtering D{\'e}bats out of the candidate
% pool whenever the classifier predicts a newspaper-only question; once
% this pre-filter is in place, the noise-suppression channel that
% TC-Chunking exploits is no longer available, and its incremental Δ
% collapses to zero. We confirm this is structural rather than a tuning
% artefact: restricting the QRe evaluation to the subset where
% TC-Chunking can still mechanically act (gold in \textit{Les D{\'e}bats}
% \textit{and} \textit{Les D{\'e}bats} retained by the router; 64\% of
% test, $n{=}283\pm3$ per seed) yields Δ Coverage@3 of $-0.002\pm0.002$,
% and increasing the strategy-selection score-$k$ from 3 to 9 leaves the
% QRe Δ within $\pm 0.001$. Table~\ref{tab:tc-chunking-by-gold-location}
% summarises the decomposition.

% \begin{table}[htbp]
% \centering
% \small
% \setlength{\tabcolsep}{5pt}
% \renewcommand{\arraystretch}{1.15}
% \begin{tabular}{l r ccc}
% \toprule
% \textbf{Subset} & $n$
%   & $\Delta$\,@3 & $\Delta$\,@5 & $\Delta$\,@10 \\
% \midrule
% Full test set                & 875        & $+5.4$  & $+6.3$  & $+6.2$  \\
% \;\;Gold in \textit{Les Débats}     & 561 (64\%) & $-0.5$  & $-0.8$  & $-0.8$  \\
% \;\;Gold not in \textit{Les Débats} & 314 (36\%) & $\mathbf{+16.2}$ & $\mathbf{+19.3}$ & $\mathbf{+19.1}$ \\
% \bottomrule
% \end{tabular}
% \caption{Naive-Retriever Δ Coverage@$k$ (pp, TC-Chunking $-$ Global
% Baseline) on the held-out test set, decomposed by whether at least one
% gold passage lies in \textit{Les Débats}. The aggregate $+5.4$~pp gain
% is carried entirely by the 36\% of questions whose gold lives in the
% two newspapers; on D{\'e}bats-gold questions TC-Chunking is mildly
% detrimental. QRe filters \textit{Les D{\'e}bats} for these same 36\%,
% absorbing the gain and reducing the QRe Δ to zero.}
% \label{tab:tc-chunking-by-gold-location}
% \end{table}

\paragraph{Gain vanishing under routing.}
Decomposing the $+5.4$~\% points Naive-Retriever Coverage@3 gain by gold
location shows it is carried entirely by the questions whose
gold lies only in the newspapers ($+16.2$~\% points), while questions with
gold in \textit{Les D\'ebats} see $-0.5$~\% points. TC-Chunking therefore helps
not by improving \textit{D\'ebats} retrieval but by making \textit{D\'ebats}
chunks less competitive in the top-$k$, freeing slots for the true
newspaper gold. QRe achieves the same outcome more directly filtering
\textit{D\'ebats} out of the candidate pool whenever the classifier
predicts a newspaper-only question which is why the TC-Chunking
$\Delta$ collapses to $\leq 0.2$~\% points once routing is active. Restricted
to the subset where TC-Chunking can still mechanically act
(gold in \textit{Les D\'ebats} \emph{and} \textit{Les D\'ebats} kept by
the router), $\Delta$ Coverage@3 is $-0.002\pm0.002$, and varying the
strategy-selection score-$k$ from 3 to 9 changes it by at most $\pm 0.001$.
Full decomposition and a residual-gain analysis are in
Appendix~\ref{sec:appendix-tc-decomposition}.


% \begin{table*}[htbp]
% \centering
% \small
% \setlength{\tabcolsep}{4pt}
% \begin{tabular}{ll cc cc}
% \toprule
% & & \multicolumn{2}{c}{\textbf{Coverage}} & \multicolumn{2}{c}{\textbf{Recall}} \\
% \cmidrule(lr){3-4} \cmidrule(lr){5-6}
% \textbf{Retrieval} & \textbf{Chunking} & \textbf{@3} & \textbf{@5} & \textbf{@3} & \textbf{@5} \\
% \midrule
% \multirow{3}{*}{Naive}
%   & Global Baseline  & 0.390 & 0.457 & 0.390 & 0.457 \\
%   & Task-Conditioned & \textbf{0.465} & \textbf{0.538} & \textbf{0.466} & \textbf{0.537} \\
%   & $\Delta$         & \cellcolor{db1}{+0.075} & \cellcolor{db1}{+0.081} & \cellcolor{db1}{+0.076} & \cellcolor{db1}{+0.080} \\
% \midrule
% \multirow{3}{*}{QRe}
%   & Global Baseline  & 0.495 & 0.578 & 0.500 & \textbf{0.580} \\
%   & Task-Conditioned & \textbf{0.511} & \textbf{0.581} & \textbf{0.511} & 0.575 \\
%   & $\Delta$         & \cellcolor{db1}{+0.016} & +0.003 & \cellcolor{db1}{+0.011} & $-$0.005 \\
% \midrule
% \multirow{3}{*}{UMS}
%   & Global Baseline  & 0.534 & 0.598 & \textbf{0.534} & 0.600 \\
%   & Task-Conditioned & \textbf{0.543} & \textbf{0.610} & \textbf{0.534} & \textbf{0.603} \\
%   & $\Delta$         & \cellcolor{db1}{+0.009} & \cellcolor{db1}{+0.012} & +0.000 & \cellcolor{db1}{+0.003} \\
% \midrule
% \multirow{3}{*}{QRe \& UMS}
%   & Global Baseline  & 0.546 & 0.607 & \textbf{0.548} & 0.610 \\
%   & Task-Conditioned & \textbf{0.555} & \textbf{0.619} & \textbf{0.548} & \textbf{0.612} \\
%   & $\Delta$         & \cellcolor{db1}{+0.009} & \cellcolor{db1}{+0.012} & +0.000 & \cellcolor{db1}{+0.002} \\
% \bottomrule
% \end{tabular}
% \caption{Coverage@$k$ and Recall@$k$ on the held-out test set ($n{=}221$):
% Global Baseline (\texttt{10000\_hierarchical}) vs.\ Task-Conditioned chunking.
% Task-conditioned gains are largest under the Naive Retriever and shrink
% substantially once routing is active.}
% \label{tab:tc-inference}
% \end{table*}


\subsection{Task-Conditioned Metadata Indexing}
\label{sec:metadata-results}
% Applying the same protocol to the orthogonal metadata axis,
% %(Section~\ref{sec:metadata-methodology})
% we hold chunking fixed at
% \texttt{10000\_hierarchical} and let each cluster pick its best metadata
% strategy on \textit{Les Débats}.
The realised gain on the held-out test set, decomposed by cluster and by
assigned strategy, is reported in
Tables~\ref{tab:metadata-wins-test-cluster} and~\ref{tab:metadata-wins-test-strategy}
(Appendix~\ref{subsec:appendix-metadata-wins}).

\paragraph{Transfer to unseen queries (RQ4).}

Table~\ref{tab:tc-inference-merged} reports
Coverage@$k$ and Recall@$k$ across the four retrievers (rows
\textit{+\,TC-Metadata}). Unlike chunking, the metadata gain persists under UMS: the gain is
largest under the Naive Retriever ($+8.9$~\% points Coverage@3, $+9.0$~\% points
Recall@3; $p<0.05$), remains substantial under UMS ($+6.0$~\% points Cov@3,
$+6.0$~\% points Rec@3; $p<0.05$) where parliamentary chunks occupy a
guaranteed retrieval quota and filtering noisy document types directly
improves that quota, vanishes under QRe ($0.0$~\% points) which already
pre-selects away from noisy sources, and is absent under
QRe\,\&\,UMS ($0.0$~\% points).
%On the held-out set only three distinct
%strategies are ever assigned in practice (\texttt{baseline}: 15 clusters,
%\texttt{exclude\_noisy}: 5, \texttt{rerank\_type\_boost}: 1).

We also evaluated the joint $14{\times}14$ chunking$\times$metadata
strategy space; the joint variant matches but never strictly improves on
TC-Metadata under any retriever, so we defer its methodology, results,
and non-significance analysis to
Appendix~\ref{sec:appendix-joint}.

% \begin{table*}[htbp]
% \centering
% \small
% \setlength{\tabcolsep}{4pt}
% \begin{tabular}{ll cc cc}
% \toprule
% & & \multicolumn{2}{c}{\textbf{Coverage}} & \multicolumn{2}{c}{\textbf{Recall}} \\
% \cmidrule(lr){3-4} \cmidrule(lr){5-6}
% \textbf{Retrieval} & \textbf{Metadata} & \textbf{@3} & \textbf{@5} & \textbf{@3} & \textbf{@5} \\
% \midrule
% \multirow{3}{*}{Naive}
%   & Global Baseline  & 0.390 & 0.457 & 0.390 & 0.457 \\
%   & Task-Conditioned & \textbf{0.498} & \textbf{0.573} & \textbf{0.502} & \textbf{0.575} \\
%   & $\Delta$         & \cellcolor{db1}{+0.108} & \cellcolor{db1}{+0.116} & \cellcolor{db1}{+0.112} & \cellcolor{db1}{+0.118} \\
% \midrule
% \multirow{3}{*}{QRe}
%   & Global Baseline  & 0.495 & 0.578 & 0.500 & 0.580 \\
%   & Task-Conditioned & \textbf{0.505} & \textbf{0.582} & \textbf{0.509} & \textbf{0.584} \\
%   & $\Delta$         & \cellcolor{db1}{+0.010} & \cellcolor{db1}{+0.004} & \cellcolor{db1}{+0.009} & \cellcolor{db1}{+0.004} \\
% \midrule
% \multirow{3}{*}{UMS}
%   & Global Baseline  & 0.534 & 0.598 & 0.534 & 0.600 \\
%   & Task-Conditioned & \textbf{0.559} & \textbf{0.607} & \textbf{0.562} & \textbf{0.610} \\
%   & $\Delta$         & \cellcolor{db1}{+0.025} & \cellcolor{db1}{+0.009} & \cellcolor{db1}{+0.028} & \cellcolor{db1}{+0.010} \\
% \midrule
% \multirow{3}{*}{QRe \& UMS}
%   & Global Baseline  & \textbf{0.546} & \textbf{0.607} & \textbf{0.548} & \textbf{0.610} \\
%   & Task-Conditioned & 0.546 & 0.607 & 0.548 & 0.610 \\
%   & $\Delta$         & +0.000 & +0.000 & +0.000 & +0.000 \\
% \bottomrule
% \end{tabular}
% \caption{Coverage@$k$ and Recall@$k$ on the held-out test set
% ($n_{\text{test}}{=}219$): Global Baseline vs.\ Task-Conditioned metadata
% indexing on \textit{Les Débats}. $k_0{=}6$ for UMS-based routings.}
% \label{tab:tc-meta-inference}
% \end{table*}

\begin{table}[htbp]
\centering
\small
\setlength{\tabcolsep}{4pt}
\renewcommand{\arraystretch}{1.1}
\resizebox{\columnwidth}{!}{%
\begin{tabular}{ll cc cc}
\toprule
\textbf{Retriever} & \textbf{Indexing strategy}
  & \multicolumn{2}{c}{\textbf{Coverage (\%)}}
  & \multicolumn{2}{c}{\textbf{Recall (\%)}} \\
\cmidrule(lr){3-4} \cmidrule(lr){5-6}
 & & @3 & @5 & @3 & @5 \\
\midrule
\multirow{3}{*}{Naive}
  & Global Baseline                 & 36.9               & 45.4               & 36.9               & 45.6               \\
  & \;+\,TC-Chunking                & 42.3$^{*}$         & 51.7$^{*}$         & 42.3$^{*}$         & 51.7$^{*}$         \\
  & \;+\,TC-Metadata                & \best{45.8}$^{*}$  & \best{55.2}$^{*}$  & \best{45.9}$^{*}$  & \best{55.3}$^{*}$  \\
\midrule
\multirow{3}{*}{QRe}
  & Global Baseline                 & 48.5               & 58.2               & 48.7               & 58.4               \\
  & \;+\,TC-Chunking                & \uline{48.5}$^{\text{ns}}$ & \uline{57.9}$^{\text{ns}}$ & \uline{48.5}$^{\text{ns}}$ & \uline{58.0}$^{\text{ns}}$ \\
  & \;+\,TC-Metadata                & \best{48.5}$^{\text{ns}}$ & \best{58.2}$^{\text{ns}}$ & \best{48.7}$^{\text{ns}}$ & \best{58.4}$^{\text{ns}}$ \\
\midrule
\multirow{3}{*}{UMS}
  & Global Baseline                 & 46.8               & 58.7               & 47.0               & 58.8               \\
  & \;+\,TC-Chunking                & 47.0$^{\text{ns}}$ & 58.5$^{\text{ns}}$ & 47.2$^{\text{ns}}$ & 58.7$^{\text{ns}}$ \\
  & \;+\,TC-Metadata                & \best{52.8}$^{*}$  & \best{63.2}$^{*}$  & \best{53.0}$^{*}$  & \best{63.4}$^{*}$  \\
\midrule
\multirow{3}{*}{QRe \& UMS}
  & Global Baseline                 & 54.7               & 64.5               & 54.9               & 64.7               \\
  & \;+\,TC-Chunking                & \uline{54.7}$^{\text{ns}}$ & \uline{64.5}$^{\text{ns}}$ & \uline{54.9}$^{\text{ns}}$ & \uline{64.7}$^{\text{ns}}$ \\
  & \;+\,TC-Metadata                & \best{54.7}$^{\text{ns}}$ & \best{64.5}$^{\text{ns}}$ & \best{54.9}$^{\text{ns}}$ & \best{64.7}$^{\text{ns}}$ \\
\bottomrule
\end{tabular}
}
\caption{Coverage@$k$ and Recall@$k$ on the held-out test set
($n{=}438$, 5 seeds), comparing each retriever's Global Baseline
(\texttt{10000\_hier.} chunking, no metadata filter, applied uniformly
to every query) to two task-conditioned variants. Under
\textit{+\,TC-Chunking}, every test question is first assigned to its
nearest question cluster, and the chunking strategy used to retrieve
its chunks is the one found optimal for that cluster on the training
split (selected from the 14 candidates of
Appendix~\ref{sec:appendix-chunking}); chunking is therefore
per-cluster, i.e.\ effectively per-question. \textit{+\,TC-Metadata}
applies the same cluster-routing logic to the metadata axis instead,
holding chunking fixed at the baseline and letting each cluster pick
its optimal metadata filter / reranker from the
14 candidates. \textbf{Bold} / \uline{underline} = best / second within block.
$^{*}p<0.05$ (Bonferroni paired $t$-test); ns = not significant.}
\label{tab:tc-inference-merged}
\end{table} in main paper
% ============================================================

\definecolor{sectiongray}{gray}{0.95}

\section{Results}
\label{sec:results}

We address four questions:
\textbf{RQ1} : does query-conditioned source selection improve retrieval over a single-index baseline?
\textbf{RQ2} : are query routing (QRe) and uniform multi-source allocation (UMS) complementary remedies for size-imbalanced retrieval?
\textbf{RQ3} : does query-conditioned chunking optimisation improve retrieval over a global baseline?
\textbf{RQ4} : does query-conditioned \emph{metadata indexing} (filters and rerankers on \textit{Les Débats}) improve retrieval over a global baseline?

\subsection{Retrieval Strategies}
\label{sec:retrieval-results}

We evaluate the four retrieval strategies of
Section~\ref{sec:retrieval-strategies} on the multi-hop subset of
HistoriQA-ThirdRepublic~\cite{historiqa-pellet}, broken down by question type:
\textit{cross-newspaper} (314 q., both gold documents from newspapers) and
\textit{newspaper $\rightarrow$ Débats} (571 q., one newspaper plus one
parliamentary document). We include comparison to BM25~\cite{bm25} for sparse retrieval and 2 graph-based external baseline: HippoRAGv2~\cite{hipporagv2} and LinearRAG~\cite{linearrag_2025}.


\begin{table*}[ht]
\centering
\small
\setlength{\tabcolsep}{4pt}
\begin{tabular*}{\textwidth}{@{\extracolsep{\fill}} l ccc ccc ccc @{}}
\toprule
 & \multicolumn{3}{c}{\textbf{Overall}} & \multicolumn{3}{c}{\textbf{Cross-newspaper} (314 q.)} & \multicolumn{3}{c}{\textbf{Newspaper-Débats} (571 q.)} \\
\cmidrule(lr){2-4} \cmidrule(lr){5-7} \cmidrule(lr){8-10}
\textbf{Configuration} & R@3 & R@5 & R@10 & R@3 & R@5 & R@10 & R@3 & R@5 & R@10 \\
\midrule
Naive Retriever~\cite{historiqa-pellet}              & 35.9 & 44.0 & 54.4 & 14.3 & 19.7 & 28.8 & 47.8 & 57.3 & 68.5 \\
BM25 (lexical)~\cite{bm25}                    & 23.1 & 29.3 & 37.8 &  4.6 &  8.3 & 13.1 & 33.3 & 40.9 & 51.4 \\
HippoRAGv2~\cite{hipporagv2}                  & 31.7 & 41.9 & 53.0 & 11.6 & 17.2 & 24.7 & 42.7 & 55.4 & 68.6 \\
LinearRAG~\cite{linearrag_2025}               & 18.2 & 27.0 & 40.1 & 6.2 & 10.4 & 20.5 & 25.1 & 36.5 & 50.9 \\
\midrule
S-RAG: QRe                                    & 48.0 & 57.0 & 68.4 & \textbf{48.4} & \textbf{56.7} & \underline{68.3} & 47.8 & 57.2 & 68.5 \\
S-RAG: UMS-Retriever                          & \underline{50.2} & \underline{57.6} & \underline{69.4} & 40.8 & 53.3 & 60.7 & \textbf{55.4} & \textbf{60.0} & \textbf{74.2} \\
S-RAG: QRe \& UMS-Retriever                   & \textbf{52.3} & \textbf{58.4} & \textbf{72.3} & \underline{46.7} & \underline{55.6} & \textbf{69.3} & \textbf{55.4} & \textbf{60.0} & \underline{73.9} \\
\bottomrule
\end{tabular*}
\caption{Recall@$k$ on multi-hop questions, with breakdown by question type.
\textbf{S-RAG}: this work. \textbf{QRe}: supervised source pre-selection.
\textbf{UMS}: equal-quota multi-source retrieval.
\textbf{BM25}: lexical baseline using the same chunking as the Naive Retriever.
\textbf{Bold} = best; \underline{underline} = second-best per column (Overall).}
\label{tab:retrieval-results}
\end{table*}

\paragraph{Overall (RQ1, RQ2).}
All three RAG configurations beat the naive, sparse retrieval and graph-based RAG baselines at every cut-off. The combined QRe\,\&\,UMS reaches Recall@3 of 52.3
($+16.4$~\% points over Naive, $+20.6\%$ over HippoRAGv2), with the gap still being high at $k{=}10$.

\paragraph{Performances by question type.}
The two strategies fix different failures. On cross-newspaper
questions, the Naive Retriever collapses (14.3 @3) as \textit{Les Débats}
floods the top-$k$ with parliamentary chunks unrelated to the query. QRe
removes this contamination and more than triples Recall@3 (48.4), while UMS alone
is less effective (40.8) because it still reserves a third of the budget
for \textit{Les Débats}. On \emph{newspaper $\leftrightarrow$ Débats} questions
the asymmetry reverses: the Naive Retriever already performs well (47.8),
QRe is neutral (the routing map keeps all three sources active), and UMS
delivers the main lift (55.4) by guaranteeing parliamentary representation
that would otherwise be crowded out.

\paragraph{Complementarity (RQ2).}
Neither routing alone nor uniform allocation alone accounts for the full
gain: each component contributes independently, and their combination is
strictly dominant overall on every metric. QRe and UMS target structurally distinct
failure modes : source contamination versus source starvation that can co-occur on a single question, which is why combining them
dominates either component in isolation.


\paragraph{Comparison to other methods.}
We include BM25~\cite{bm25} as a sparse-retrieval baseline using the same chunking
as the Naive Retriever. BM25 underperforms the dense baselines on every
cut-off and on both question types, with gap on cross-newspaper questions of 4.6 vs.\ 14.3 Recall@3 against the
Naive Retriever. Lexical matching is particularly ill-suited to this
benchmark: questions are paraphrased rather than quoted from the
articles, and the OCR'd 19th-century newspaper text introduces surface
variation (orthographic and typographic) that dense embeddings absorb
but term-frequency matching does not. We therefore retain dense
retrieval as the underlying index for all Structured retriever variants.

Despite its graph-based index, HippoRAG v2 trails the Naive Retriever on cross-newspaper questions (11.6 vs.\ 14.3 @3) and underperforms all structured retriever variants across every metric. LinearRAG~\cite{linearrag_2025} is the weakest baseline
overall (Recall@3 of 18.2 vs.\ 35.9 for the Naive Retriever) and degrades
sharpest on cross-newspaper questions (6.2 @3, 10.4 @5, 20.5 @10). Even on the
more favourable \emph{newspaper $\leftrightarrow$ D\'ebats} split it remains
below the Naive Retriever at every cut-off (25.1 / 36.5 / 50.9 vs.\ 47.8 / 57.3 / 68.5). Those results suggests that graphs are poorly constructed and that propagation breaks down when evidence is spread across heterogeneous corpora marked by OCR errors and inconsistent chunk sizes.

\subsection{Task/Query-Conditioned Chunking}
\label{sec:indexing-results}

We now analyse whether per-cluster selection of a chunking strategy improves
retrieval over a single global best. The comparison of the 14 candidate
strategies under the four retrievers is in Appendix~\ref{subsec:appendix-global-strategy-comparison}
(Table~\ref{tab:recall-coverage-k3}).
The realised gain on the held-out test set, decomposed by cluster and by
assigned strategy, is reported in
Tables~\ref{tab:chunking-wins-test-cluster} and~\ref{tab:chunking-wins-summary}
(Appendix~\ref{subsec:appendix-cluster-wins}), strategy-level qualitative wins are discussed in
Appendix~\ref{sec:strategy-wins}. A non-trivial subset of clusters
benefits from a strictly different strategy than the global baseline,
motivating the held-out transfer evaluation below.

\paragraph{Transfer to unseen queries (RQ3).}
% The oracle analysis shows the potential for strong retrieval gains when
% the appropriate chunking strategy is chosen, but it assumes the best
% strategy is known in advance. We test the realised gain using the
% inference procedure described in Section~\ref{sec:indexing-inference}.
Results are reported in Table~\ref{tab:tc-inference-merged} (rows
\textit{+\,TC-Chunking}) alongside the metadata-axis variant
discussed below. Under the Naive Retriever the inference pipeline realises
$+5.4$~\% points Coverage@3 and $+6.1$~\% points Recall@5 ($p<0.05$), confirming
that a single low-cost classification step at query time delivers
meaningful gain. Gains are negative but non significant when Qre is active (at most $-0.4$ points) and are negligible and non-significant under UMS retriever (at most $+0.2$~\% points, $p>0.05$ for all routing
configurations), indicating that task-conditioned chunking and source
routing address entirely overlapping failure modes: the gains chunking
provides for the naive retriever are fully recovered by routing alone.

% \paragraph{Mechanism: where TC-Chunking actually helps.}
% A post-hoc decomposition of the Naive-Retriever gain by gold location
% clarifies why TC-Chunking and QRe overlap. Splitting the 875 test
% questions by whether at least one gold passage lies in \textit{Les
% D{\'e}bats} (the long-form parliamentary source, $n{=}561$, 64\%) or
% only in the two newspapers ($n{=}314$, 36\%), the $+5.4$~pp average
% Coverage@3 gain decomposes into $-0.5$~pp on gold-in-D{\'e}bats
% questions and $+16.2$~pp on the remaining 36\%. The benefit of
% TC-Chunking is therefore not better \textit{D{\'e}bats} retrieval but
% \textit{less aggressive competition} from D{\'e}bats chunks in the
% top-$k$, freeing slots for the true newspaper gold. QRe achieves the
% same outcome more directly by filtering D{\'e}bats out of the candidate
% pool whenever the classifier predicts a newspaper-only question; once
% this pre-filter is in place, the noise-suppression channel that
% TC-Chunking exploits is no longer available, and its incremental Δ
% collapses to zero. We confirm this is structural rather than a tuning
% artefact: restricting the QRe evaluation to the subset where
% TC-Chunking can still mechanically act (gold in \textit{Les D{\'e}bats}
% \textit{and} \textit{Les D{\'e}bats} retained by the router; 64\% of
% test, $n{=}283\pm3$ per seed) yields Δ Coverage@3 of $-0.002\pm0.002$,
% and increasing the strategy-selection score-$k$ from 3 to 9 leaves the
% QRe Δ within $\pm 0.001$. Table~\ref{tab:tc-chunking-by-gold-location}
% summarises the decomposition.

% \begin{table}[htbp]
% \centering
% \small
% \setlength{\tabcolsep}{5pt}
% \renewcommand{\arraystretch}{1.15}
% \begin{tabular}{l r ccc}
% \toprule
% \textbf{Subset} & $n$
%   & $\Delta$\,@3 & $\Delta$\,@5 & $\Delta$\,@10 \\
% \midrule
% Full test set                & 875        & $+5.4$  & $+6.3$  & $+6.2$  \\
% \;\;Gold in \textit{Les Débats}     & 561 (64\%) & $-0.5$  & $-0.8$  & $-0.8$  \\
% \;\;Gold not in \textit{Les Débats} & 314 (36\%) & $\mathbf{+16.2}$ & $\mathbf{+19.3}$ & $\mathbf{+19.1}$ \\
% \bottomrule
% \end{tabular}
% \caption{Naive-Retriever Δ Coverage@$k$ (pp, TC-Chunking $-$ Global
% Baseline) on the held-out test set, decomposed by whether at least one
% gold passage lies in \textit{Les Débats}. The aggregate $+5.4$~pp gain
% is carried entirely by the 36\% of questions whose gold lives in the
% two newspapers; on D{\'e}bats-gold questions TC-Chunking is mildly
% detrimental. QRe filters \textit{Les D{\'e}bats} for these same 36\%,
% absorbing the gain and reducing the QRe Δ to zero.}
% \label{tab:tc-chunking-by-gold-location}
% \end{table}

\paragraph{Gain vanishing under routing.}
Decomposing the $+5.4$~\% points Naive-Retriever Coverage@3 gain by gold
location shows it is carried entirely by the questions whose
gold lies only in the newspapers ($+16.2$~\% points), while questions with
gold in \textit{Les D\'ebats} see $-0.5$~\% points. TC-Chunking therefore helps
not by improving \textit{D\'ebats} retrieval but by making \textit{D\'ebats}
chunks less competitive in the top-$k$, freeing slots for the true
newspaper gold. QRe achieves the same outcome more directly filtering
\textit{D\'ebats} out of the candidate pool whenever the classifier
predicts a newspaper-only question which is why the TC-Chunking
$\Delta$ collapses to $\leq 0.2$~\% points once routing is active. Restricted
to the subset where TC-Chunking can still mechanically act
(gold in \textit{Les D\'ebats} \emph{and} \textit{Les D\'ebats} kept by
the router), $\Delta$ Coverage@3 is $-0.002\pm0.002$, and varying the
strategy-selection score-$k$ from 3 to 9 changes it by at most $\pm 0.001$.
Full decomposition and a residual-gain analysis are in
Appendix~\ref{sec:appendix-tc-decomposition}.


% \begin{table*}[htbp]
% \centering
% \small
% \setlength{\tabcolsep}{4pt}
% \begin{tabular}{ll cc cc}
% \toprule
% & & \multicolumn{2}{c}{\textbf{Coverage}} & \multicolumn{2}{c}{\textbf{Recall}} \\
% \cmidrule(lr){3-4} \cmidrule(lr){5-6}
% \textbf{Retrieval} & \textbf{Chunking} & \textbf{@3} & \textbf{@5} & \textbf{@3} & \textbf{@5} \\
% \midrule
% \multirow{3}{*}{Naive}
%   & Global Baseline  & 0.390 & 0.457 & 0.390 & 0.457 \\
%   & Task-Conditioned & \textbf{0.465} & \textbf{0.538} & \textbf{0.466} & \textbf{0.537} \\
%   & $\Delta$         & \cellcolor{db1}{+0.075} & \cellcolor{db1}{+0.081} & \cellcolor{db1}{+0.076} & \cellcolor{db1}{+0.080} \\
% \midrule
% \multirow{3}{*}{QRe}
%   & Global Baseline  & 0.495 & 0.578 & 0.500 & \textbf{0.580} \\
%   & Task-Conditioned & \textbf{0.511} & \textbf{0.581} & \textbf{0.511} & 0.575 \\
%   & $\Delta$         & \cellcolor{db1}{+0.016} & +0.003 & \cellcolor{db1}{+0.011} & $-$0.005 \\
% \midrule
% \multirow{3}{*}{UMS}
%   & Global Baseline  & 0.534 & 0.598 & \textbf{0.534} & 0.600 \\
%   & Task-Conditioned & \textbf{0.543} & \textbf{0.610} & \textbf{0.534} & \textbf{0.603} \\
%   & $\Delta$         & \cellcolor{db1}{+0.009} & \cellcolor{db1}{+0.012} & +0.000 & \cellcolor{db1}{+0.003} \\
% \midrule
% \multirow{3}{*}{QRe \& UMS}
%   & Global Baseline  & 0.546 & 0.607 & \textbf{0.548} & 0.610 \\
%   & Task-Conditioned & \textbf{0.555} & \textbf{0.619} & \textbf{0.548} & \textbf{0.612} \\
%   & $\Delta$         & \cellcolor{db1}{+0.009} & \cellcolor{db1}{+0.012} & +0.000 & \cellcolor{db1}{+0.002} \\
% \bottomrule
% \end{tabular}
% \caption{Coverage@$k$ and Recall@$k$ on the held-out test set ($n{=}221$):
% Global Baseline (\texttt{10000\_hierarchical}) vs.\ Task-Conditioned chunking.
% Task-conditioned gains are largest under the Naive Retriever and shrink
% substantially once routing is active.}
% \label{tab:tc-inference}
% \end{table*}


\subsection{Task-Conditioned Metadata Indexing}
\label{sec:metadata-results}
% Applying the same protocol to the orthogonal metadata axis,
% %(Section~\ref{sec:metadata-methodology})
% we hold chunking fixed at
% \texttt{10000\_hierarchical} and let each cluster pick its best metadata
% strategy on \textit{Les Débats}.
The realised gain on the held-out test set, decomposed by cluster and by
assigned strategy, is reported in
Tables~\ref{tab:metadata-wins-test-cluster} and~\ref{tab:metadata-wins-test-strategy}
(Appendix~\ref{subsec:appendix-metadata-wins}).

\paragraph{Transfer to unseen queries (RQ4).}

Table~\ref{tab:tc-inference-merged} reports
Coverage@$k$ and Recall@$k$ across the four retrievers (rows
\textit{+\,TC-Metadata}). Unlike chunking, the metadata gain persists under UMS: the gain is
largest under the Naive Retriever ($+8.9$~\% points Coverage@3, $+9.0$~\% points
Recall@3; $p<0.05$), remains substantial under UMS ($+6.0$~\% points Cov@3,
$+6.0$~\% points Rec@3; $p<0.05$) where parliamentary chunks occupy a
guaranteed retrieval quota and filtering noisy document types directly
improves that quota, vanishes under QRe ($0.0$~\% points) which already
pre-selects away from noisy sources, and is absent under
QRe\,\&\,UMS ($0.0$~\% points).
%On the held-out set only three distinct
%strategies are ever assigned in practice (\texttt{baseline}: 15 clusters,
%\texttt{exclude\_noisy}: 5, \texttt{rerank\_type\_boost}: 1).

We also evaluated the joint $14{\times}14$ chunking$\times$metadata
strategy space; the joint variant matches but never strictly improves on
TC-Metadata under any retriever, so we defer its methodology, results,
and non-significance analysis to
Appendix~\ref{sec:appendix-joint}.

% \begin{table*}[htbp]
% \centering
% \small
% \setlength{\tabcolsep}{4pt}
% \begin{tabular}{ll cc cc}
% \toprule
% & & \multicolumn{2}{c}{\textbf{Coverage}} & \multicolumn{2}{c}{\textbf{Recall}} \\
% \cmidrule(lr){3-4} \cmidrule(lr){5-6}
% \textbf{Retrieval} & \textbf{Metadata} & \textbf{@3} & \textbf{@5} & \textbf{@3} & \textbf{@5} \\
% \midrule
% \multirow{3}{*}{Naive}
%   & Global Baseline  & 0.390 & 0.457 & 0.390 & 0.457 \\
%   & Task-Conditioned & \textbf{0.498} & \textbf{0.573} & \textbf{0.502} & \textbf{0.575} \\
%   & $\Delta$         & \cellcolor{db1}{+0.108} & \cellcolor{db1}{+0.116} & \cellcolor{db1}{+0.112} & \cellcolor{db1}{+0.118} \\
% \midrule
% \multirow{3}{*}{QRe}
%   & Global Baseline  & 0.495 & 0.578 & 0.500 & 0.580 \\
%   & Task-Conditioned & \textbf{0.505} & \textbf{0.582} & \textbf{0.509} & \textbf{0.584} \\
%   & $\Delta$         & \cellcolor{db1}{+0.010} & \cellcolor{db1}{+0.004} & \cellcolor{db1}{+0.009} & \cellcolor{db1}{+0.004} \\
% \midrule
% \multirow{3}{*}{UMS}
%   & Global Baseline  & 0.534 & 0.598 & 0.534 & 0.600 \\
%   & Task-Conditioned & \textbf{0.559} & \textbf{0.607} & \textbf{0.562} & \textbf{0.610} \\
%   & $\Delta$         & \cellcolor{db1}{+0.025} & \cellcolor{db1}{+0.009} & \cellcolor{db1}{+0.028} & \cellcolor{db1}{+0.010} \\
% \midrule
% \multirow{3}{*}{QRe \& UMS}
%   & Global Baseline  & \textbf{0.546} & \textbf{0.607} & \textbf{0.548} & \textbf{0.610} \\
%   & Task-Conditioned & 0.546 & 0.607 & 0.548 & 0.610 \\
%   & $\Delta$         & +0.000 & +0.000 & +0.000 & +0.000 \\
% \bottomrule
% \end{tabular}
% \caption{Coverage@$k$ and Recall@$k$ on the held-out test set
% ($n_{\text{test}}{=}219$): Global Baseline vs.\ Task-Conditioned metadata
% indexing on \textit{Les Débats}. $k_0{=}6$ for UMS-based routings.}
% \label{tab:tc-meta-inference}
% \end{table*}

\begin{table}[htbp]
\centering
\small
\setlength{\tabcolsep}{4pt}
\renewcommand{\arraystretch}{1.1}
\resizebox{\columnwidth}{!}{%
\begin{tabular}{ll cc cc}
\toprule
\textbf{Retriever} & \textbf{Indexing strategy}
  & \multicolumn{2}{c}{\textbf{Coverage (\%)}}
  & \multicolumn{2}{c}{\textbf{Recall (\%)}} \\
\cmidrule(lr){3-4} \cmidrule(lr){5-6}
 & & @3 & @5 & @3 & @5 \\
\midrule
\multirow{3}{*}{Naive}
  & Global Baseline                 & 36.9               & 45.4               & 36.9               & 45.6               \\
  & \;+\,TC-Chunking                & 42.3$^{*}$         & 51.7$^{*}$         & 42.3$^{*}$         & 51.7$^{*}$         \\
  & \;+\,TC-Metadata                & \best{45.8}$^{*}$  & \best{55.2}$^{*}$  & \best{45.9}$^{*}$  & \best{55.3}$^{*}$  \\
\midrule
\multirow{3}{*}{QRe}
  & Global Baseline                 & 48.5               & 58.2               & 48.7               & 58.4               \\
  & \;+\,TC-Chunking                & \uline{48.5}$^{\text{ns}}$ & \uline{57.9}$^{\text{ns}}$ & \uline{48.5}$^{\text{ns}}$ & \uline{58.0}$^{\text{ns}}$ \\
  & \;+\,TC-Metadata                & \best{48.5}$^{\text{ns}}$ & \best{58.2}$^{\text{ns}}$ & \best{48.7}$^{\text{ns}}$ & \best{58.4}$^{\text{ns}}$ \\
\midrule
\multirow{3}{*}{UMS}
  & Global Baseline                 & 46.8               & 58.7               & 47.0               & 58.8               \\
  & \;+\,TC-Chunking                & 47.0$^{\text{ns}}$ & 58.5$^{\text{ns}}$ & 47.2$^{\text{ns}}$ & 58.7$^{\text{ns}}$ \\
  & \;+\,TC-Metadata                & \best{52.8}$^{*}$  & \best{63.2}$^{*}$  & \best{53.0}$^{*}$  & \best{63.4}$^{*}$  \\
\midrule
\multirow{3}{*}{QRe \& UMS}
  & Global Baseline                 & 54.7               & 64.5               & 54.9               & 64.7               \\
  & \;+\,TC-Chunking                & \uline{54.7}$^{\text{ns}}$ & \uline{64.5}$^{\text{ns}}$ & \uline{54.9}$^{\text{ns}}$ & \uline{64.7}$^{\text{ns}}$ \\
  & \;+\,TC-Metadata                & \best{54.7}$^{\text{ns}}$ & \best{64.5}$^{\text{ns}}$ & \best{54.9}$^{\text{ns}}$ & \best{64.7}$^{\text{ns}}$ \\
\bottomrule
\end{tabular}
}
\caption{Coverage@$k$ and Recall@$k$ on the held-out test set
($n{=}438$, 5 seeds), comparing each retriever's Global Baseline
(\texttt{10000\_hier.} chunking, no metadata filter, applied uniformly
to every query) to two task-conditioned variants. Under
\textit{+\,TC-Chunking}, every test question is first assigned to its
nearest question cluster, and the chunking strategy used to retrieve
its chunks is the one found optimal for that cluster on the training
split (selected from the 14 candidates of
Appendix~\ref{sec:appendix-chunking}); chunking is therefore
per-cluster, i.e.\ effectively per-question. \textit{+\,TC-Metadata}
applies the same cluster-routing logic to the metadata axis instead,
holding chunking fixed at the baseline and letting each cluster pick
its optimal metadata filter / reranker from the
14 candidates. \textbf{Bold} / \uline{underline} = best / second within block.
$^{*}p<0.05$ (Bonferroni paired $t$-test); ns = not significant.}
\label{tab:tc-inference-merged}
\end{table}